% =====================================================================
% draft_sections/framework_instantiations.tex
% Subsection: how MRMA drops into each of the three backbones.
% Intended placement: end of Section~\ref{sec:method} (after the MRMA module
% and Algorithm~\ref{alg:mrma}), before Experiments.
% Uses \textsc{...} spellouts; if the macros in framework_intro.tex are added
% to the preamble they may be substituted throughout.
% =====================================================================

\subsection{Backbone Instantiations}
\label{sec:instantiations}

MRMA (Section~\ref{sec:method}) is defined purely as an operation on a
channel-mixing gate: it takes the pre-softmax logit matrix that couples the $N$
channel tokens and adds the fixed relation gate $\mathbf{G}^{(r)}$ of
Eq.~(\ref{eq:mrma}) on a per-head basis.
Any backbone that exposes such an $N{\times}N$ channel-coupling matrix can host
it \emph{unchanged}: the three mask buffers, the cyclic head-to-relation
assignment $r_i{=}i\bmod 3$, and the graph precomputation of
Algorithm~\ref{alg:precompute} are identical across hosts.
We verify this by instantiating MRMA in three backbones from three distinct
architecture families.
In every case the mask buffers are the \emph{only} addition, so each
instantiation has exactly the parameter count of its own backbone---the
comparison is like-for-like.

\textbf{\mr{} ${=}$ \itrans{} ${+}$ MRMA (inverted Transformer).}
This is our primary instantiation and the one developed in
Section~\ref{sec:method}.
\itrans{}~\cite{liu2023itransformer} embeds each channel as a token and applies
multi-head self-attention across the $N$ channel tokens; MRMA replaces that
self-attention with the masked variant of Eq.~(\ref{eq:mrma}).
Token embedding, feed-forward sublayers, layer normalisation, RevIN, and the
output projection are untouched.
\emph{What changes:} the pre-softmax channel-attention logits, gated per head.
\emph{What stays:} everything else, at identical parameter count.

\textbf{\textsc{SECrossformer} ${=}$ \textsc{CrossFormer} ${+}$ MRMA
(cross-dimension Transformer).}
\textsc{CrossFormer}~\cite{zhang2023crossformer} mixes channels through a
\emph{learned} two-stage cross-dimension attention (a router-augmented attention
over dimension segments) rather than a single flat channel self-attention.
MRMA drops in at exactly the stage where channels are coupled: the same three
relation gates are added to the cross-dimension attention logits before their
softmax, restricting each head to its assigned relation neighbourhood.
The two-stage router, the patch/segment embedding, and CrossFormer's
hierarchical encoder are left in place.
This instantiation is the most informative stress test of transferability,
because CrossFormer already \emph{learns} a channel graph through its two-stage
router: MRMA here constrains an already-structured mixing rather than a flat
global attention. As reported in Section~\ref{sec:frameworkresults}, this is
exactly where the mechanism's benefit is smallest---the static masks help in only
5 of 16 cells---which is the informative negative half of the transfer story.
\emph{What changes:} the cross-dimension attention logits are masked.
\emph{What stays:} the router, segment embedding, and hierarchy---zero added
parameters.

\textbf{\textsc{StructMamba} ${=}$ GRU-based state-space forecaster ${+}$ MRMA
(recurrent state-space).}
The third backbone is not a Transformer at all: it is a GRU-based state-space
forecaster in the selective-SSM forecasting family
(cf.~S-Mamba~\cite{wang2024smamba}, DMamba~\cite{chen2026dmamba}), which mixes
channels through a gated recurrent state-update rather than an explicit
attention softmax.
MRMA enters as a binary gate on the channel-coupling term of the state update:
a channel pair $(i,j)$ blocked by all of its assigned relation masks does not
contribute to the recurrent mix, exactly as a $-\infty$ logit removes it from a
softmax.
The GRU recurrence, the state-space parameterisation, and the projection heads
are unchanged, and the fair backbone for this instantiation is precisely
\emph{the same forecaster with the masks set to all-ones} (no relational
prior).
On this recurrent backbone the masks deliver their largest per-cell gains:
against the all-ones control, error drops in all 16 PEMS cells (mean $+13.8\%$,
up to $+26.6\%$; Section~\ref{sec:frameworkresults}). That a non-attention,
recurrent host benefits at least as much as the attention host is the strongest
evidence that the relational graphs themselves, not the attention softmax they
were first formulated in, carry the benefit.
\emph{What changes:} the channel-coupling term of the recurrent state update is
gated by the masks.
\emph{What stays:} the GRU state-space core and all learned weights---again zero
added parameters.

In all three cases the precomputed graphs are the identical $(3{\times}N{\times}N)$
buffer from Algorithm~\ref{alg:precompute}; no host-specific graph, retraining of
the masks, or architecture-specific tuning of the relation construction is
required.
The framework thus reduces to a single reusable artefact---the mask array---and
a two-line change at each backbone's channel-mixing step.
